% The Weierstrass wp-Function  (source §7; added: full convergence proof,
% rigorous periodicity/evenness, Eisenstein series and Laurent expansion,
% complete derivation of wp'^2 = 4wp^3 - g2 wp - g3, the roots e_i, the
% discriminant Delta != 0, wp as a degree-2 branched cover, and generation of
% the field of elliptic functions)
\chapter{The Weierstrass \texorpdfstring{$\wp$}{wp}-Function}
\label{ch:weierstrass_wp}

Riemann--Roch told us (\cref{ch:riemann_roch}) that the torus must carry a
function with a double pole at \(O\); it did not produce one. This chapter
constructs it explicitly --- the Weierstrass function \(\wp\) --- proves the
series converges, and derives by hand the relation
\(\wp'^2 = 4\wp^3 - g_2\wp - g_3\) predicted in \cref{rem:rr-cubic-preview}. The
roots of the right-hand side are distinct (\(\Delta \ne 0\)), so the cubic is
smooth; \(\wp\) realizes the torus as a double cover of \(\CP^1\); and the pair
\((\wp, \wp')\) generates the whole field of elliptic functions. Everything the
later chapters need about \(\wp\) is established here.

\section{Definition and Convergence}
\label{sec:wp-def}

\begin{definition}
\label{def:wp}
For a lattice \(\Lambda\), the \emph{Weierstrass function} is
\[
  \wp(z) = \frac{1}{z^2}
  + \sum_{\lambda \in \Lambda \setminus \{0\}}
    \left( \frac{1}{(z-\lambda)^2} - \frac{1}{\lambda^2} \right).
\]
\end{definition}

The subtracted \(1/\lambda^2\) is what makes the series converge; without it the
sum of \(1/(z-\lambda)^2\) would diverge. Convergence rests on a counting
estimate for the lattice.

\begin{lemma}[Lattice sum estimate]
\label{lem:wp-latsum}
The number of lattice points with \(|\lambda| \le R\) is \(O(R^2)\), and
consequently
\[
  \sum_{\lambda \in \Lambda \setminus \{0\}} \frac{1}{|\lambda|^{s}} < \infty
  \qquad\text{for every } s > 2 .
\]
\end{lemma}

\begin{proof}
By \cref{ex:lat-discrete} the lattice points are separated, so disjoint disks of
a fixed radius \(\rho\) about them lie within \(|\lambda| \le R + \rho\); comparing
areas, their number is at most \(\pi(R+\rho)^2/(\pi\rho^2) = O(R^2)\). Grouping
the nonzero lattice points into annuli \(n \le |\lambda| < n+1\), each annulus
contains \(O(n)\) points, each contributing \(\le n^{-s}\), so
\(\sum_\lambda |\lambda|^{-s} \le \sum_{n\ge 1} O(n)\,n^{-s}
= \sum_{n \ge 1} O(n^{1-s})\), which converges for \(s > 2\).
\end{proof}

\begin{proposition}
\label{prop:wp-converges}
The series defining \(\wp\) converges absolutely and uniformly on every compact
subset of \(\C \setminus \Lambda\). Hence \(\wp\) is meromorphic on \(\C\), with a
double pole and zero residue at each lattice point and no other poles.
\end{proposition}

\begin{proof}
Fix \(R > 0\) and consider \(|z| \le R\). For \(|\lambda| \ge 2R\),
\[
  \left| \frac{1}{(z-\lambda)^2} - \frac{1}{\lambda^2} \right|
  = \left| \frac{z(2\lambda - z)}{\lambda^2 (z-\lambda)^2} \right|
  \le \frac{R\,(2|\lambda| + R)}{|\lambda|^2 (|\lambda| - R)^2}
  \le \frac{R \cdot 3|\lambda|}{|\lambda|^2 (|\lambda|/2)^2}
  = \frac{12 R}{|\lambda|^3}.
\]
By \cref{lem:wp-latsum} with \(s = 3\), \(\sum_\lambda |\lambda|^{-3} < \infty\),
so by the Weierstrass \(M\)-test the tail \(\sum_{|\lambda| \ge 2R}\) converges
uniformly on \(|z| \le R\); the finitely many remaining terms are meromorphic
there. Thus \(\wp\) is meromorphic, and near a lattice point \(\lambda_0\) the
term \(1/(z-\lambda_0)^2\) supplies a double pole with zero residue while the rest
is holomorphic.
\end{proof}

\section{Periodicity and Evenness}
\label{sec:wp-periodic}

\begin{proposition}
\label{prop:wp-even-periodic}
\(\wp\) is even, \(\wp(-z) = \wp(z)\), and \(\Lambda\)-periodic,
\(\wp(z+\lambda) = \wp(z)\). Its derivative
\[
  \wp'(z) = -2 \sum_{\lambda \in \Lambda} \frac{1}{(z-\lambda)^3}
\]
is \(\Lambda\)-periodic and odd, with a triple pole at each lattice point.
\end{proposition}

\begin{proof}
\emph{Evenness.} Replacing \(z\) by \(-z\) and reindexing \(\lambda \mapsto
-\lambda\) (which permutes \(\Lambda\)),
\(\wp(-z) = \tfrac1{z^2} + \sum_\lambda\bigl(\tfrac1{(z+\lambda)^2} -
\tfrac1{\lambda^2}\bigr) = \wp(z)\).

\emph{The derivative.} Differentiating the uniformly convergent series term by
term gives \(\wp'(z) = -2\sum_{\lambda}(z-\lambda)^{-3}\), where now the sum runs
over \emph{all} \(\lambda \in \Lambda\) (the \(\lambda = 0\) term is the
derivative of \(1/z^2\)). This sum is manifestly periodic: for \(\mu \in \Lambda\),
substituting \(z \mapsto z+\mu\) permutes the summands \(\lambda \mapsto
\lambda - \mu\), so \(\wp'(z+\mu) = \wp'(z)\). It is odd because
\(\wp = \) even implies \(\wp'\) odd (or directly from the symmetric sum).

\emph{Periodicity of \(\wp\).} For a basis vector \(\omega\), the function
\(g(z) = \wp(z+\omega) - \wp(z)\) has \(g'(z) = \wp'(z+\omega) - \wp'(z) = 0\), so
\(g\) is a constant \(c\). Evaluating at \(z = -\omega/2\) (not a lattice point)
and using evenness, \(c = \wp(\omega/2) - \wp(-\omega/2) = 0\). Hence
\(\wp(z+\omega) = \wp(z)\) for each basis vector, so \(\wp\) is
\(\Lambda\)-periodic.
\end{proof}

Thus \(\wp\) is an elliptic function of order two (one double pole per period),
the minimal order allowed by \cref{cor:ell-order2}, and \(\wp'\) is elliptic of
order three.

\section{Eisenstein Series and the Laurent Expansion}
\label{sec:wp-eisenstein}

\begin{definition}
\label{def:wp-eisenstein}
For an even integer \(m \ge 4\) the \emph{Eisenstein series} of weight \(m\) is
\[
  G_m = G_m(\Lambda) = \sum_{\lambda \in \Lambda \setminus \{0\}} \frac{1}{\lambda^{m}},
\]
which converges by \cref{lem:wp-latsum}. For odd \(m\) the sum vanishes by the
symmetry \(\lambda \mapsto -\lambda\).
\end{definition}

\begin{proposition}
\label{prop:wp-laurent}
Near \(z = 0\),
\[
  \wp(z) = \frac{1}{z^2} + \sum_{k \ge 1} (2k+1)\,G_{2k+2}\, z^{2k}
  = \frac{1}{z^2} + 3 G_4 z^2 + 5 G_6 z^4 + O(z^6).
\]
\end{proposition}

\begin{proof}
For \(0 < |z| < |\lambda|\), expand \(\frac{1}{(z-\lambda)^2} =
\frac{1}{\lambda^2}\bigl(1 - z/\lambda\bigr)^{-2}
= \frac{1}{\lambda^2} \sum_{n \ge 0} (n+1)\,(z/\lambda)^n\), so
\[
  \frac{1}{(z-\lambda)^2} - \frac{1}{\lambda^2}
  = \sum_{n \ge 1} (n+1)\,\frac{z^n}{\lambda^{n+2}} .
\]
Summing over \(\lambda \ne 0\) (absolute convergence justifies interchange) and
adding \(1/z^2\),
\[
  \wp(z) = \frac{1}{z^2} + \sum_{n \ge 1} (n+1)\,G_{n+2}\,z^n .
\]
The odd-\(n\) terms drop (\(G_{n+2} = 0\) for \(n\) odd), leaving the stated even
series with \(a_2 = 3G_4\), \(a_4 = 5G_6\).
\end{proof}

\section{The Differential Equation}
\label{sec:wp-ode}

\begin{theorem}[Weierstrass differential equation]
\label{thm:wp-ode}
With \(g_2 = 60 G_4\) and \(g_3 = 140 G_6\),
\[
  \boxed{\ \wp'^2 = 4\wp^3 - g_2 \wp - g_3 .\ }
\]
\end{theorem}

\begin{proof}
From \cref{prop:wp-laurent}, \(\wp = z^{-2} + 3G_4 z^2 + 5G_6 z^4 + O(z^6)\) and
\(\wp' = -2z^{-3} + 6 G_4 z + 20 G_6 z^3 + O(z^5)\). Squaring and cubing,
\[
  \wp'^2 = 4 z^{-6} - 24 G_4 z^{-2} - 80 G_6 + O(z^2),
  \qquad
  4\wp^3 = 4 z^{-6} + 36 G_4 z^{-2} + 60 G_6 + O(z^2).
\]
Subtracting, the order-six poles cancel and
\[
  \wp'^2 - 4\wp^3 = -60 G_4 z^{-2} - 140 G_6 + O(z^2).
\]
Adding \(g_2 \wp = 60 G_4 z^{-2} + O(z^2)\) and the constant \(g_3 = 140 G_6\)
cancels the remaining negative power and constant:
\[
  F := \wp'^2 - 4\wp^3 + g_2 \wp + g_3 = O(z^2) .
\]
Now \(F\) is elliptic (a polynomial in \(\wp, \wp'\)), and its only possible poles
are at lattice points; the computation shows \(F\) is holomorphic at \(0\), and by
periodicity at every lattice point, so \(F\) is holomorphic on the compact torus,
hence constant (\cref{thm:ell-const}). Since \(F = O(z^2) \to 0\) as
\(z \to 0\), that constant is \(0\). Therefore \(\wp'^2 = 4\wp^3 - g_2\wp - g_3\).
\end{proof}

\begin{keyidea}{The birth of the cubic}
The equation \(\wp'^2 = 4\wp^3 - g_2\wp - g_3\) is not guessed. A specific
polynomial in \(\wp\) and \(\wp'\) was arranged to have no pole at the origin;
being elliptic and pole-free it is constant, and the Laurent expansion shows the
constant is zero. The map \(z \mapsto (\wp(z), \wp'(z))\) therefore lands on the
plane cubic \(y^2 = 4x^3 - g_2 x - g_3\) --- the embedding of
\cref{ch:embedding}.
\end{keyidea}

\section{The Roots, the Discriminant, and the Double Cover}
\label{sec:wp-roots}

The cubic factors over \(\C\) as \(4\wp^3 - g_2\wp - g_3 = 4(\wp - e_1)(\wp - e_2)
(\wp - e_3)\). The roots have a direct meaning on the torus.

\begin{proposition}
\label{prop:wp-roots}
Let \(\tfrac{\omega_1}{2}, \tfrac{\omega_2}{2}, \tfrac{\omega_1+\omega_2}{2}\) be
the three nonzero half-periods and set
\(e_1 = \wp(\tfrac{\omega_1}{2})\), \(e_2 = \wp(\tfrac{\omega_2}{2})\),
\(e_3 = \wp(\tfrac{\omega_1+\omega_2}{2})\). Then:
\begin{enumerate}
\item \(\wp'\) vanishes exactly at the three half-periods, each a simple zero;
\item \(4\wp^3 - g_2\wp - g_3 = 4(\wp-e_1)(\wp-e_2)(\wp-e_3)\);
\item \(e_1, e_2, e_3\) are distinct.
\end{enumerate}
\end{proposition}

\begin{proof}
(1) \(\wp'\) is odd and periodic, so at a half-period \(\omega/2\), writing
\(-\omega/2 \equiv \omega/2 \pmod\Lambda\), \(\wp'(\omega/2) = \wp'(-\omega/2) =
-\wp'(\omega/2)\), hence \(\wp'(\omega/2) = 0\). That is three zeros; \(\wp'\) has
order three (triple pole at \(O\)), so by \cref{thm:ell-valence} it has exactly
three zeros, and these are simple.

(2) By \cref{thm:wp-ode}, \(\wp'^2 = 4\wp^3 - g_2\wp - g_3\) vanishes wherever
\(\wp' = 0\), i.e.\ at the half-periods, where \(\wp\) takes the values
\(e_1, e_2, e_3\). So \(e_1, e_2, e_3\) are the roots of the cubic
\(4x^3 - g_2 x - g_3\), giving the factorization.

(3) Suppose \(e_1 = e_2\). The function \(\wp - e_1\) has order two (double pole
at \(O\)), so exactly two zeros. It vanishes at \(\omega_1/2\), and there
\((\wp - e_1)' = \wp' = 0\), so that zero is double --- using up both zeros. Then
\(\wp - e_1\) cannot also vanish at \(\omega_2/2\), contradicting
\(\wp(\omega_2/2) = e_2 = e_1\). Hence \(e_1 \ne e_2\), and similarly all three
are distinct.
\end{proof}

\begin{definition}
\label{def:wp-discriminant}
The \emph{discriminant} of \(\Lambda\) is \(\Delta = g_2^3 - 27 g_3^2\).
\end{definition}

\begin{corollary}
\label{cor:wp-smooth}
\(\Delta = 16\,(e_1 - e_2)^2 (e_2 - e_3)^2 (e_1 - e_3)^2 \ne 0\). Equivalently the
cubic \(y^2 = 4x^3 - g_2 x - g_3\) is nonsingular.
\end{corollary}

\begin{proof}
For \(4x^3 - g_2 x - g_3 = 4(x-e_1)(x-e_2)(x-e_3)\) the elementary symmetric
relations give \(e_1+e_2+e_3 = 0\), and a standard discriminant computation yields
\(g_2^3 - 27 g_3^2 = 16 \prod_{i<j}(e_i - e_j)^2\). By \cref{prop:wp-roots}(3) the
\(e_i\) are distinct, so \(\Delta \ne 0\). A plane cubic \(y^2 = p(x)\) is singular
exactly when \(p\) has a repeated root, which \(\Delta \ne 0\) excludes.
\end{proof}

The geometric content is that \(\wp\) presents the torus as a branched double
cover of the sphere (\cref{fig:branchedcover}): a generic value is attained twice
(\(\wp\) has order two), while the four values \(e_1, e_2, e_3, \infty\) are
attained once, at the half-periods and at \(O\). The Riemann--Hurwitz count
\(\chi(T) = 2\,\chi(\CP^1) - (\text{branch points})\) reads
\(0 = 2 \cdot 2 - 4\), confirming genus one.

\begin{figure}[ht]
\centering
\begin{tikzpicture}[scale=1.0, every node/.style={font=\small}]
  % Top: precise torus with the four 2-torsion (half-lattice) points marked.
  \begin{scope}[yshift=2.0cm, scale=0.85]
    \drawprecisetorus
    \fill[formrust!85!black] (\torusMarkhalfsumx,\torusMarkhalfsumy) circle (3pt);
    \node[formrust!60!black, above, font=\scriptsize]
      at (\torusMarkhalfsumx,\torusMarkhalfsumy) {$[\tfrac{\omega_1+\omega_2}{2}]$};
    \fill[formrust!85!black] (\torusMarkzerox,\torusMarkzeroy) circle (3pt);
    \node[formrust!60!black, below right, font=\scriptsize]
      at (\torusMarkzerox,\torusMarkzeroy) {$[\tfrac{\omega_1}{2}]$};
    \fill[formrust!85!black] (\torusMarkhalftwox,\torusMarkhalftwoy) circle (3pt);
    \node[formrust!60!black, below left, font=\scriptsize]
      at (\torusMarkhalftwox,\torusMarkhalftwoy) {$[\tfrac{\omega_2}{2}]$};
    \fill[formrust!85!black] (\torusMarkhalfonex,\torusMarkhalfoney) circle (3pt);
    \node[formrust!60!black, right, font=\scriptsize]
      at (\torusMarkhalfonex,\torusMarkhalfoney) {$[0]$};
    \node[formteal!55!black] at (-3.0,1.0) {$T$};
  \end{scope}
  \draw[-{Latex[length=3mm]}, formplum, line width=1.1pt] (0,0.7) -- (0,-0.1);
  \node[formplum] at (0.35,0.3) {$\wp$, deg 2};
  \begin{scope}[yshift=-1.5cm, scale=1.1]
    \drawprecisesphere
    \fill[formrust!85!black] (\sphereNorthx,\sphereNorthy) circle (2.4pt);
    \node[formrust!60!black, above, font=\scriptsize]
      at (\sphereNorthx,\sphereNorthy) {$\infty$};
    \fill[formrust!85!black] (-0.707,-0.354) circle (2.4pt);
    \node[formrust!60!black, left, font=\scriptsize] at (-0.78,-0.354) {$e_1$};
    \fill[formrust!85!black] (0,-0.500) circle (2.4pt);
    \node[formrust!60!black, below, font=\scriptsize] at (0,-0.56) {$e_2$};
    \fill[formrust!85!black] (0.707,-0.354) circle (2.4pt);
    \node[formrust!60!black, right, font=\scriptsize] at (0.78,-0.354) {$e_3$};
    \node[formteal!55!black] at (-1.65,0.6) {$\CP^1$};
  \end{scope}
  \node[draw=black!40, rounded corners=2pt, fill=formivory, align=left, anchor=west,
        inner sep=4pt] at (3.4,0.4)
    {The map $\wp:T\to\CP^1$ is a\\$2:1$ branched cover\\
     with $4$ branch values:\\
     three roots $e_1,e_2,e_3$ and $\infty$.\\
     Riemann--Hurwitz:\\
     $0=2{\cdot}2-4{\cdot}(2-1)$.};
\end{tikzpicture}
\caption{\(\wp\) as a degree-two branched cover \(T \to \CP^1\). The four branch
points are the half-periods (mapping to \(e_1, e_2, e_3\)) and \(O\) (mapping to
\(\infty\)). Torus and sphere are rendered from \texttt{sagemath/08\_torus\_render.sage}
and \texttt{sagemath/10\_sphere\_render.sage}.}
\label{fig:branchedcover}
\end{figure}

\section{Generation of the Field of Elliptic Functions}
\label{sec:wp-generation}

We promised in \cref{ch:elliptic_liouville} that \(\wp, \wp'\) generate all
elliptic functions. We can now prove it, and identify the field structure.

\begin{theorem}
\label{thm:wp-generation}
The field of \(\Lambda\)-elliptic functions is \(\C(\wp, \wp')\). Moreover
\(\wp' \notin \C(\wp)\), so \(\C(\wp,\wp')\) is a degree-two extension of
\(\C(\wp)\), with \(\wp'\) satisfying the minimal relation
\(\wp'^2 = 4\wp^3 - g_2\wp - g_3\).
\end{theorem}

\begin{proof}
\emph{Minimality.} Every element of \(\C(\wp)\) is even, since \(\wp\) is even, so
a rational function \(R(\wp)\) satisfies \(R(\wp)(-z) = R(\wp)(z)\). But \(\wp'\)
is odd and not identically zero, so \(\wp' \notin \C(\wp)\). As \(\wp'^2 \in
\C(\wp)\) by \cref{thm:wp-ode}, the extension \(\C(\wp,\wp')/\C(\wp)\) has degree
exactly two.

\emph{Generation.} Let \(f\) be elliptic. Split it into even and odd parts,
\[
  f = f_+ + f_-, \quad
  f_+(z) = \tfrac12\bigl(f(z) + f(-z)\bigr), \quad
  f_-(z) = \tfrac12\bigl(f(z) - f(-z)\bigr),
\]
both elliptic. Then \(f_-/\wp'\) is even (odd over odd), so it suffices to show
every even elliptic function lies in \(\C(\wp)\). Let \(h\) be even elliptic.
Since \(h(-z) = h(z)\), its divisor is invariant under \(z \mapsto -z\), so its
zeros and poles occur in pairs \(\{[a], [-a]\}\) of equal order --- except at the
four two-torsion points, where \([a] = [-a]\). At a two-torsion point \(a\) the
order of \(h\) is \emph{even}: in the local coordinate \(w = z - a\) the
involution \(z \mapsto -z\) is \(w \mapsto -w\) (because \(2a \in \Lambda\)), and
\(h(a+w) = h(a-w)\), so \(h\) is an even function of \(w\) and vanishes (or blows
up) to even order.

Now match \(\divisor(h)\) with a rational function of \(\wp\). For a
non-two-torsion pair where \(h\) has order \(m\), the factor
\((\wp - \wp(a))^{m}\) has divisor \(m\bigl([a] + [-a] - 2[O]\bigr)\). For a
two-torsion point \(a\) where \(h\) has order \(2m\), the factor
\((\wp - \wp(a))^{m}\) has divisor \(2m\bigl([a] - [O]\bigr)\), because
\(\wp - \wp(a)\) has a \emph{double} zero at the half-period \(a\) (there
\(\wp' = 0\), \cref{prop:wp-roots}). Let \(R(\wp)\) be the product of these
factors, the exponent positive at zeros of \(h\) and negative at its poles. Then
\(\divisor(h)\) and \(\divisor R(\wp)\) agree at every point except possibly
\(O\); both have degree \(0\) (\cref{thm:div-degzero}), so their coefficients at
\(O\) agree too, and \(\divisor(h) = \divisor R(\wp)\). Hence \(h/R(\wp)\) is
elliptic with no zeros or poles, so it is constant, and \(h \in \C(\wp)\).
Therefore \(f = f_+ + (f_-/\wp')\,\wp' \in \C(\wp,\wp')\).
\end{proof}

\section{The \texorpdfstring{$\sigma$}{sigma}-Function and the Converse of Abel}
\label{sec:wp-sigma}

The converse of Abel's law --- that a degree-zero divisor with Abel sum zero is the
divisor of an elliptic function --- was deferred in
\cref{ch:divisors,ch:riemann_roch} to an existence theorem. We now discharge it
\emph{constructively}, building the function explicitly from a third Weierstrass
function.

\begin{definition}
\label{def:wp-sigma}
The Weierstrass \emph{sigma function} is
\[
  \sigma(z) = z \prod_{\lambda \in \Lambda\setminus\{0\}}
  \Bigl(1 - \tfrac{z}{\lambda}\Bigr)
  \exp\!\Bigl(\tfrac{z}{\lambda} + \tfrac12\bigl(\tfrac{z}{\lambda}\bigr)^2\Bigr).
\]
\end{definition}

\begin{proposition}
\label{prop:wp-sigma}
The product converges normally on compact sets, so \(\sigma\) is entire; it is odd,
\(\sigma(-z) = -\sigma(z)\), and its zeros are exactly the lattice points, each
simple.
\end{proposition}

\begin{proof}
Taking logarithms, the factor at \(\lambda\) contributes
\[
  \log\!\Bigl(1 - \tfrac{z}{\lambda}\Bigr) + \tfrac{z}{\lambda}
  + \tfrac12\bigl(\tfrac{z}{\lambda}\bigr)^2
  = -\sum_{m\ge3}\frac1m\Bigl(\tfrac{z}{\lambda}\Bigr)^m
  = O\bigl(|\lambda|^{-3}\bigr)
\]
for \(|z|\le R\), \(|\lambda|\ge 2R\); by \cref{lem:wp-latsum} with \(s = 3\) the
sum converges normally, so the product converges to an entire function whose zeros
are those of
its factors --- \(z = 0\) from the leading \(z\), and \(z = \lambda\) from
\(1 - z/\lambda\) --- each simple. For oddness, re-index \(\lambda \mapsto -\lambda\):
\(\sigma(-z) = -z\prod(1 + z/\lambda)\exp(-z/\lambda + \tfrac12(z/\lambda)^2)
= -\sigma(z)\).
\end{proof}

The logarithmic derivative \(\zeta = \sigma'/\sigma\) is the Weierstrass
\emph{zeta function}
\[
  \zeta(z) = \frac1z + \sum_{\lambda\neq0}
  \Bigl(\frac{1}{z-\lambda} + \frac1\lambda + \frac{z}{\lambda^2}\Bigr),
\]
odd, with a simple pole of residue \(1\) at each lattice point. Differentiating
term by term, \(\zeta'(z) = -\tfrac1{z^2} + \sum_{\lambda\neq0}\bigl(-\tfrac1{(z-\lambda)^2}
+ \tfrac1{\lambda^2}\bigr) = -\wp(z)\) (\cref{def:wp}); equivalently
\(\wp = -(\log\sigma)''\), so this \(\sigma\) is the one attached to the book's
\(\wp\).

\begin{proposition}[Quasi-periodicity]
\label{prop:wp-sigma-quasi}
For each basis period \(\omega_k\) there is a constant \(\eta_k = 2\zeta(\omega_k/2)\)
with
\[
  \zeta(z+\omega_k) = \zeta(z) + \eta_k,
  \qquad
  \sigma(z+\omega_k) = -\,\sigma(z)\,\exp\!\bigl(\eta_k(z + \tfrac{\omega_k}{2})\bigr).
\]
\end{proposition}

\begin{proof}
Since \(\zeta' = -\wp\) is \(\Lambda\)-periodic (\cref{prop:wp-even-periodic}),
\(\zeta(z+\omega_k) - \zeta(z)\) is constant; evaluating at \(z = -\omega_k/2\) and
using \(\zeta\) odd gives the constant \(2\zeta(\omega_k/2) =: \eta_k\). As
\((\log\sigma)' = \zeta\), \(\tfrac{d}{dz}\log\tfrac{\sigma(z+\omega_k)}{\sigma(z)}
= \eta_k\), so \(\sigma(z+\omega_k) = \sigma(z)\,e^{\eta_k z + c_k}\) for a constant
\(c_k\). Setting \(z = -\omega_k/2\) and using \(\sigma(-\omega_k/2) =
-\sigma(\omega_k/2) \neq 0\) gives \(e^{c_k - \eta_k\omega_k/2} = -1\), which is the
stated formula.
\end{proof}

\begin{remark}
\label{rem:wp-legendre}
The quasi-periods obey the Legendre relation \(\eta_1\omega_2 - \eta_2\omega_1 =
2\pi i\) for a positively oriented basis (integrate \(\zeta\) around a period
parallelogram); we will not need it.
\end{remark}

\begin{theorem}[Converse of Abel, constructively]
\label{thm:wp-converse}
Let \(D = \sum_{i=1}^n [a_i] - \sum_{i=1}^n [b_i]\) be a divisor of degree zero on
\(T\) with Abel sum zero, \(\sum_i a_i \equiv \sum_i b_i \pmod\Lambda\). Then \(D\)
is principal, \(D = \divisor(F)\), with
\[
  F(z) = \frac{\prod_{i=1}^n \sigma(z - a_i)}{\prod_{i=1}^n \sigma(z - b_i)} .
\]
\end{theorem}

\begin{proof}
Adjusting one representative by a lattice vector --- which changes neither its class
nor \(D\) --- we may assume \(\sum_i a_i = \sum_i b_i\) in \(\C\). By
\cref{prop:wp-sigma} the \(a_i\) are simple zeros and the \(b_i\) simple poles of
\(F\) (with multiplicity if repeated), so \(\divisor(F) = D\). For periodicity, by
\cref{prop:wp-sigma-quasi} each basis period gives
\[
  \frac{F(z+\omega_k)}{F(z)}
  = \frac{\prod_i\bigl(-e^{\eta_k(z - a_i + \omega_k/2)}\bigr)}
         {\prod_i\bigl(-e^{\eta_k(z - b_i + \omega_k/2)}\bigr)}
  = \exp\!\Bigl(\eta_k\bigl(\textstyle\sum_i b_i - \sum_i a_i\bigr)\Bigr) = 1 :
\]
the \((-1)^n\) factors and the \(nz + n\omega_k/2\) terms cancel because the counts
are equal (degree zero), and the exponent vanishes because \(\sum a_i = \sum b_i\)
(Abel sum zero). So \(F\) is elliptic with \(\divisor(F) = D\), and \(D\) is
principal.
\end{proof}

\begin{remark}
\label{rem:wp-selfcontained}
\Cref{thm:wp-converse} is the converse of Abel's law; it completes the isomorphism
\(\Pic^0(T) \cong T\) (\cref{thm:div-pic-iso,cor:rr-degree-one}) with no appeal to
Riemann--Roch. Together with the elementary counts of \cref{prop:rr-counts}, every
load-bearing result in the book is now proved from within it. Riemann--Roch
(\cref{thm:rr}) remains as the general theorem, illuminating but unneeded.
\end{remark}

\section{Summary}
\label{sec:wp-summary}

The Weierstrass function \(\wp(z) = z^{-2} + \sum_{\lambda \ne 0}\bigl((z-\lambda)^{-2}
- \lambda^{-2}\bigr)\) converges (\cref{prop:wp-converges}) to an even elliptic
function of order two, with odd derivative \(\wp'\) of order three
(\cref{prop:wp-even-periodic}). Its Laurent expansion has Eisenstein coefficients
\(G_4, G_6\) (\cref{prop:wp-laurent}), and a pole-cancellation argument yields the
differential equation \(\wp'^2 = 4\wp^3 - g_2\wp - g_3\) with \(g_2 = 60G_4\),
\(g_3 = 140G_6\) (\cref{thm:wp-ode}). The roots \(e_i = \wp(\omega_i/2)\) are the
values of \(\wp\) at the half-periods and are distinct, so the discriminant
\(\Delta = g_2^3 - 27g_3^2\) is nonzero and the cubic is smooth
(\cref{prop:wp-roots,cor:wp-smooth}); \(\wp\) exhibits \(T\) as a double cover of
\(\CP^1\) branched over \(e_1, e_2, e_3, \infty\). Finally \(\C(\wp, \wp')\) is the
full field of elliptic functions, a quadratic extension of \(\C(\wp)\)
(\cref{thm:wp-generation}). Finally the sigma function \(\sigma\)
(\cref{def:wp-sigma}), a third Weierstrass function, builds an elliptic function
with any prescribed balanced divisor, proving the converse of Abel's law
constructively (\cref{thm:wp-converse}) --- so \(\Pic^0(T) \cong T\) needs no
Riemann--Roch. The next chapter turns \((\wp, \wp')\) into the embedding of
\cref{thm:main}.

\section*{Exercises}
\addcontentsline{toc}{section}{Exercises}

\begin{exercise}
\label{ex:wp-convergence}
Prove \cref{lem:wp-latsum} in detail, and use it to justify term-by-term
differentiation of the \(\wp\)-series, giving
\(\wp'(z) = -2\sum_{\lambda}(z-\lambda)^{-3}\).
\end{exercise}

\begin{exercise}
\label{ex:wp-g2g3}
From the Laurent expansion derive \(a_2 = 3G_4\) and \(a_4 = 5G_6\), and confirm
\(g_2 = 20a_2 = 60G_4\) and \(g_3 = 28a_4 = 140G_6\). (Match the
\(z^{-2}\) and constant terms of \(\wp'^2 - 4\wp^3\).)
\end{exercise}

\begin{exercise}
\label{ex:wp-roots-distinct}
Give the discriminant computation \(g_2^3 - 27g_3^2 = 16\prod_{i<j}(e_i-e_j)^2\)
starting from \(e_1+e_2+e_3 = 0\), \(e_1e_2+e_1e_3+e_2e_3 = -g_2/4\),
\(e_1e_2e_3 = g_3/4\). Conclude \(\Delta \ne 0\) from
\cref{prop:wp-roots}(3).
\end{exercise}

\begin{exercise}
\label{ex:wp-field}
Show every elliptic function \(f\) can be written \(f = A(\wp) + B(\wp)\,\wp'\)
with \(A, B\) rational, by splitting \(f\) into even and odd parts as in
\cref{thm:wp-generation}. Where is \(\wp' \notin \C(\wp)\) used?
\end{exercise}

\begin{exercise}
\label{ex:wp-recurrence}
Differentiate \(\wp'^2 = 4\wp^3 - g_2\wp - g_3\) to obtain
\(\wp'' = 6\wp^2 - \tfrac12 g_2\). Expanding both sides in Laurent series, derive
a recurrence expressing each Eisenstein coefficient \(G_{2k}\) (\(k \ge 4\)) as a
polynomial in \(G_4\) and \(G_6\). (Thus all \(G_{2k}\) are determined by
\(G_4, G_6\).)
\end{exercise}

\begin{exercise}
\label{ex:wp-sigma-quasi}
Complete the proof of \cref{prop:wp-sigma-quasi}: starting from \(\zeta(z+\omega_k)
= \zeta(z) + \eta_k\) and \((\log\sigma)' = \zeta\), integrate to get
\(\sigma(z+\omega_k) = \sigma(z)\,e^{\eta_k z + c_k}\), then determine the constant
\(c_k\) by setting \(z = -\omega_k/2\) and using that \(\sigma\) is odd, obtaining
\(\sigma(z+\omega_k) = -\sigma(z)\,e^{\eta_k(z + \omega_k/2)}\).
\end{exercise}

\begin{exercise}
\label{ex:wp-converse-check}
In the proof of \cref{thm:wp-converse}, verify the cancellation
\(F(z+\omega_k)/F(z) = \exp\!\bigl(\eta_k(\sum_i b_i - \sum_i a_i)\bigr)\) in
detail, and identify exactly where the two hypotheses are used: that the number of
zeros equals the number of poles (degree zero), and that \(\sum_i a_i = \sum_i b_i\)
(Abel sum zero). Which hypothesis fails, and what goes wrong, for the divisor
\([a] - [0]\) with \(a \notin \Lambda\)?
\end{exercise}
